\section*{Synthèse}
\addcontentsline{toc}{section}{Synthèse}
\markboth{Synthèse}{}

Ce projet reprend les notebooks d'un projet antérieur de modélisation énergétique
(prédiction de la consommation et des émissions du parc immobilier de Seattle) et
les porte en production : une application en ligne, un modèle versionné, une
chaîne d'intégration continue et un journal de décisions.

L'audit de l'existant a produit le constat structurant du projet. Le modèle
initial affichait un $R^2$ de $0{,}81$ sur les émissions, mais il consommait des
variables calculées \emph{à partir de la consommation réelle} : un utilisateur
capable de les renseigner connaît déjà la réponse qu'il vient chercher. Retirer
cette fuite coûte dix points de performance mesurée et rend le modèle utilisable
sur un bâtiment jamais relevé — c'est-à-dire sur le cas d'usage.

\begin{constat}{Le résultat le plus solide du projet n'est pas une performance, c'est une limite mesurée.}
Parmi 163 immeubles de bureaux de taille comparable, la consommation au pied
carré varie d'un facteur $3{,}4$. Aucun modèle utilisant ces variables ne peut
départager ces bâtiments. Un intervalle construit sur cette seule dispersion
irréductible vaudrait un facteur $2{,}24$ ; le modèle livré produit un facteur
$2{,}03$. \textbf{Il est donc déjà plus serré que l'écart naturel entre bâtiments
que rien ne distingue.} L'incertitude affichée n'est pas un défaut de
modélisation, c'est une propriété du problème.
\end{constat}

Chaque estimation est donc accompagnée d'une fourchette, produite par
\emph{prédiction conforme}, dont la couverture est garantie sans hypothèse sur la
loi des erreurs. L'interface annonce explicitement que l'outil sert à classer un
parc et à prioriser des audits, pas à chiffrer un bâtiment.

Trois arbitrages ont été tranchés par la mesure plutôt que par l'opinion : un
modèle de fondation écarté malgré le meilleur score du banc d'essai, une
plateforme d'hébergement changée en cours de route sous contrainte commerciale
externe, et un dispositif de supervision volontairement non implémenté. Les trois
sont documentés avec leurs chiffres, y compris ceux qui contredisent l'intuition
de départ.

\vspace{0.6em}
\noindent\textbf{Livrables} : application publique, dépôt de code, 61 tests
automatisés, chaîne CI/CD déployante, registre d'expérimentations, et un journal
de seize écarts documentés vis-à-vis de la solution auditée.

\newpage
